\documentclass[11pt]{article}
\usepackage[T1]{fontenc}
\usepackage{lmodern,amsmath,amssymb,amsthm,mathtools,microtype}
\usepackage[margin=1in]{geometry}
\usepackage[hidelinks]{hyperref}
\newtheorem{theorem}{Proposed theorem}
\newtheorem{lemma}{Lemma}
\newtheorem{proposition}{Proposition}
\newtheorem{remark}{Remark}
\newcommand{\ReH}{\operatorname{Re}_{H}}
\newcommand{\ImH}{\operatorname{Im}_{H}}
\newcommand{\norm}[1]{\left\lVert#1\right\rVert}
\newcommand{\abs}[1]{\left|#1\right|}
\newcommand{\ip}[2]{\left\langle#1,#2\right\rangle}
\newcommand{\dd}{\,\mathrm d}
\newcommand{\ex}{e_x}
\newcommand{\ey}{e_y}
\newcommand{\ee}{\mathrm e}
\title{Uniform absorption asymmetry on a fixed frequency band\\
\large A constructive draft for a strict extinction floor}
\author{Research working note}
\date{17 October 2026}
\begin{document}
\maketitle
\begin{abstract}
We give a proposed constructive extension of the compact two-layer
prescribed-extinction argument to a fixed frequency band. For every band
bounded away from zero, every positive center-frequency extinction, and every
strictly smaller positive extinction floor, the argument yields compact
reciprocal nonmagnetic dielectric bodies whose worst-case normalized absorption
asymmetry approaches one. Each body has a single passive causal constitutive
law, made from finitely many Lorentz oscillators, and obeys arbitrary prescribed
upper dielectric and loss caps throughout the band. The additional ingredient
is a passive matrix-valued dispersion construction that shapes polarization
conversion; a fixed-axis weak plate alone gives a weaker floor range.
The argument is an analytical research draft. It is separate from the existing
manuscript and frozen formalization. The boundary case in which the floor equals
the prescribed center extinction is not established.
\end{abstract}

\section{Statement and exact scope}
Fix
\[
 B=[\omega_1,\omega_2],\qquad
 0<\omega_1<\omega_c<\omega_2,\qquad
 k=\omega/c_{\rm v},
\]
where $c_{\rm v}$ is the vacuum speed. Use the opposite incident electric fields
$f_\pm=\ey\ee^{\pm ikz}$ with equal incident intensity. The object comprises two
homogeneous layers with common square aperture $S_R=(-R/2,R/2)^2$, thicknesses
$h,d>0$, and vacuum outside. Each layer has a spatially constant, symmetric
$3\times3$ dielectric tensor depending on frequency; $\mu=I$. The total
extinction, absorption and full-sphere scattering cross sections satisfy
\[
 X_+=X_-=X=A_\pm+S_\pm,\qquad
 \Gamma=\frac{\abs{A_+-A_-}}{X}.
\]
These are plane-wave cross sections, not fractions of a finite incident power.

\begin{theorem}[Fixed band with a strict extinction floor]\label{thm:main}
Fix $P>1$, $Q>0$, $E_*>0$, and $0<E_{\min}<E_*$. For every
$0<\epsilon<1$, the construction below supplies a finite two-layer body and
one passive causal reciprocal dielectric law in each layer such that, for
every $\omega\in B$,
\begin{gather}
 I\preceq\ReH\varepsilon(\omega)\preceq PI,\qquad
 0\prec\ImH\varepsilon(\omega)\preceq QI,\label{eq:caps}\\
 X(\omega_c)=E_*,\qquad
 X(\omega)\ge E_{\min},\qquad
 1-\epsilon<\Gamma(\omega)<1.\label{eq:target}
\end{gather}
Each susceptibility is a finite sum of damped Lorentz oscillators with real
positive semidefinite matrix residues, and tends to zero at infinite frequency.
Consequently, within this resource-unrestricted class,
\[
 \sup_{\text{admissible bodies}}\ \inf_{\omega\in B}\Gamma(\omega)=1.
\]
The supremum is not attained.
\end{theorem}

Material caps are imposed on $B$, not on all real frequencies. There is no
common aperture, thickness, volume, oscillator-strength, number-of-oscillators
or positive-loss-floor bound across the sequence. Frequencies do not select
different bodies: all geometrical and oscillator parameters are fixed for each
selected body. The constituent oscillator axes are fixed; the principal axes
of their summed tensor may vary with frequency. This is permitted by the
two-layer tensor class of the base manuscript \cite{base}.

The proof uses the weak-Maxwell volume representation, radiation identity and
optical theorem established in the base manuscript. The new frequency-uniform
localization, matrix factorization, dispersion shaping and parameter selection
are supplied here. This is not a new Lean verification or an independent audit
of all inherited Maxwell foundations.

\section{Passive dispersion can shape the conversion amplitude}
For $F>0$, $\Omega>0$, $\gamma>0$, define
\[
 \ell(z)=\frac{F}{\Omega^2-z^2-i\gamma z}.
\]
Its poles lie in the lower half-plane, it obeys the real-time symmetry
$\ell(-\overline z)=\overline{\ell(z)}$, and for $\Im z>0$,
\begin{equation}\label{eq:herglotz}
 \Im\{z\ell(z)\}
 =
 F\frac{(\Im z)(\Omega^2+\abs z^2)+\gamma\abs z^2}
        {\abs{\Omega^2-z^2-i\gamma z}^2}>0.
\end{equation}
Thus $z\ell(z)$ is Herglotz, the passivity condition for electric susceptibility.
Multiplication by a real positive semidefinite matrix, finite summation, and
congruence by a fixed real matrix preserve this condition. If
$\Omega>\omega_2$, both real and imaginary parts of $\ell(\omega)$ are positive
throughout $B$. Matrix oscillator positivity is also central in general
broadband scattering representations \cite{oscillators}; the local construction
used here is explicit.

\begin{lemma}[Signed resolvent approximation]\label{lem:density}
The real linear span of $(t-x)^{-1}$, $t>T>\omega_2^2$, is dense in
$C([\omega_1^2,\omega_2^2];\mathbb R)$.
\end{lemma}
\begin{proof}
An annihilating finite signed measure $\mu$ would have Stieltjes transform
$\int(t-x)^{-1}\dd\mu(x)=0$ for every $t>T$. Expanding at infinity gives
$\int x^n\dd\mu(x)=0$ for every nonnegative integer $n$. Polynomial density on
the compact interval implies $\mu=0$. The Hahn--Banach theorem and the Riesz
representation theorem now give density. Every approximant is a finite sum;
repeated poles may be merged.
\end{proof}

Fix $a_0>0$ and choose
\begin{equation}\label{eq:tau}
 \tau=\frac{2\pi}{\omega_1+\omega_2}.
\end{equation}
In particular $0<\tau\omega_1<\tau\omega_2<2\pi$. The positive continuous target
\[
 f(\omega)=\frac{a_0}{2\sin(\tau\omega/2)}
\]
is bounded on $B$. Lemma~\ref{lem:density} gives a finite real rational sum
\[
 r(\omega)=\sum_{j=1}^m\frac{\beta_j}{t_j-\omega^2}
\]
arbitrarily close to $f$. The $\beta_j$ may have either sign. For any
$\eta_j>0$, the residue
\[
 M_j=
 \begin{pmatrix}\eta_j&\beta_j\\
                 \beta_j&\beta_j^2/\eta_j\end{pmatrix}
\]
is positive semidefinite. Choose the $\eta_j$ small enough that
$\sum_j\eta_j/(t_j-\omega_2^2)$ is arbitrarily small. Define the lossless
intermediate matrix
\begin{equation}\label{eq:C0}
 C_0(\omega)=
 \frac{a_0T_0}{T_0-\omega^2}\ex\ex^T+
 \sum_{j=1}^m\frac{M_j}{t_j-\omega^2}+
 \frac{\kappa I_2}{T_*-\omega^2}
 =
 \begin{pmatrix}a(\omega)&b(\omega)\\ b(\omega)&c(\omega)\end{pmatrix}.
\end{equation}
All poles exceed $\omega_2^2$. Take $T_0$ sufficiently large and $\kappa>0$
sufficiently small. Then $C_0\succ0$, $a$ is arbitrarily close to $a_0$, and
$b=r$ is arbitrarily close to $f$. The entry $c$ may be extremely large, but
it is finite and continuous. Its unrestricted size is an essential resource
in this construction.

Replace every denominator in \eqref{eq:C0} by
$t-\omega^2-i\gamma_t\omega$, using positive damping also for the isotropic
term. The resulting $C_\gamma$ obeys
\[
 \ReH C_\gamma\succ0,\qquad \ImH C_\gamma\succ0,\qquad
 C_\gamma\longrightarrow C_0
\]
uniformly on $B$ as all dampings tend to zero. It is a finite passive reciprocal
oscillator sum with high-frequency decay. All these choices describe one
material law across the whole band.

\section{The weak thick plate and the selective absorber}
Set $D_\theta=\operatorname{diag}(1,\theta)$, $0\le\theta\le1$. Before final
isotropic regularization, use the plate susceptibility
\begin{equation}\label{eq:plate}
 \chi_{W,N,\theta}(\omega)=
 \frac1N D_\theta C_\gamma(\omega)D_\theta
 \quad\text{in the $xy$ block},\qquad
 d_N=\frac{2c_{\rm v}\tau N}{a_0},
\end{equation}
with zero $z$ susceptibility. The absorber has susceptibility
\begin{equation}\label{eq:absorber}
 \chi_{A,N}(z)=\alpha_{A,N}(z)\ex\ex^T,\quad
 \alpha_{A,N}(z)=\frac{4\omega_2^2/N}
                         {4\omega_2^2-z^2-i\omega_2 z},\quad
 h_N=\frac{N^2c_{\rm v}}{\omega_c}.
\end{equation}
For each fixed oscillator matrix, sufficiently large $N$ gives strict upper
margins in both material caps, uniformly in $\theta\in[0,1]$ and $\omega\in B$.
The lower real bound follows from positive real oscillator coefficients and
positive semidefinite residues.

\subsection{Uniform absorber limit}
Write $x=\omega/\omega_2$ and $a_B=\omega_1/\omega_2>0$. Then
$\alpha_{A,N}=4/[N(4-x^2-ix)]$. On $B$,
\[
 \abs{\alpha_{A,N}}\le\frac4{3N},\qquad
 \Im\alpha_{A,N}\ge\frac{a_B}{4N}.
\]
For the passive square root $n_A=\sqrt{1+\alpha_{A,N}}$,
$\Im n_A\ge c_A/N$ with $c_A>0$ independent of $\omega$.
The slab propagation amplitude $p_A=\exp(ikn_Ah_N)$ therefore satisfies
$\abs{p_A}\le\exp(-cN)$. Its Fresnel amplitude is $O(N^{-1})$. The exact
scalar slab formulas give, uniformly on $B$,
\begin{equation}\label{eq:absorber-limit}
 r_A=O(N^{-1}),\qquad t_A=O(\exp(-cN)).
\end{equation}
The $y$ channel remains vacuum-equivalent in this intermediate absorber.

\subsection{Uniform plate limit after normalization}
First take $C_\gamma=C_0$, and put
$M_\theta=D_\theta C_0D_\theta$ and $t=\tau\omega/a_0$. The index matrix is
$n_N=(I+M_\theta/N)^{1/2}$. All its slab amplitudes are matrix functions of
$M_\theta$:
\begin{align}
 r_0&=(I-n_N)(I+n_N)^{-1},& p_N&=\exp(ikd_Nn_N),\nonumber\\
 R_W&=r_0(I-p_N^2)(I-r_0^2p_N^2)^{-1},&
 T_W&=(I-r_0^2)p_N(I-r_0^2p_N^2)^{-1}.\label{eq:matrix-slab}
\end{align}
For fixed $C_0$, these expressions are holomorphic in $\theta$ on a fixed
complex disk for all sufficiently large $N$. There,
\[
 n_N=I+\frac{M_\theta}{2N}+O(N^{-2}),\quad
 R_W=O(N^{-1}),\quad
 \frac{T_W}{v_W}=Q_\theta+O(N^{-1}),\quad
 Q_\theta=\exp(itM_\theta),\quad v_W=\ee^{ikd_N}.
\]
The estimates are uniform in real frequency on $B$ and on the complex
$\theta$ disk. Although complex $\theta$ is only an analytic device, its
propagation matrices have a finite uniform norm on the disk; hence the
denominators in \eqref{eq:matrix-slab} are invertible for large $N$.

Parity under $\theta\mapsto-\theta$ makes offdiagonal entries odd and diagonal
entries even. At $\theta=0$ the $yy$ slab is exactly vacuum. Applying Cauchy
estimates on a smaller disk therefore gives
\begin{gather}
 (R_W)_{xy}/\theta=O(N^{-1}),\qquad
 (R_W)_{yy}/\theta^2=O(N^{-1}),\nonumber\\
 (T_W/v_W-Q_\theta)_{xy}/\theta=O(N^{-1}),\qquad
 (T_W/v_W-Q_\theta)_{yy}/\theta^2=O(N^{-1}).\label{eq:normalized-slab}
\end{gather}
These are the needed normalized estimates; unnormalized convergence alone
would not control vanishing coupling.

Direct expansion of the exponential, with $\phi=ta$, gives
\begin{align}
 (Q_\theta)_{xy}
 &=\theta\frac ba(\ee^{i\phi}-1)+O(\theta^3),\label{eq:Qxy}\\
 (Q_\theta)_{yy}
 &=1+\theta^2\left[
 itc+\frac{b^2}{a^2}(\ee^{i\phi}-1-i\phi)\right]+O(\theta^4).\label{eq:Qyy}
\end{align}
The large $c$ contributes only an imaginary quadratic forward phase and does
not enter the quadratic extinction coefficient.

\subsection{Actual stack amplitudes and the normalized spectrum}
Let $P_x=\ex\ex^T$, $P_y=\ey\ey^T$, $v_A=\ee^{ikh_N}$,
$R_A=r_AP_x$, $T_A=t_AP_x+v_AP_y$. Exact cascade matching gives
\begin{align*}
 t_+&=T_W(I-R_AR_W)^{-1}T_A\ey,\\
 r_+&=T_AR_W(I-R_AR_W)^{-1}T_A\ey,\\
 t_-&=T_A(I-R_WR_A)^{-1}T_W\ey,\\
 r_-&=\left[R_W+T_WR_A(I-R_WR_A)^{-1}T_W\right]\ey.
\end{align*}
The unlisted direct $R_A\ey$ term is zero. Reciprocal transmission gives the
same co-polar forward coefficient. With $v=v_Av_W$, the planar extinction is
\[
 X_\infty=2\Re(1-\ey^Tt_\pm/v),
\]
and scattered transmission is $t_\pm-v\ey$. Define
$Y_\infty=X_\infty/\theta^2$, and similarly normalized scattering
$C_\infty^\pm$ and absorption $B_\infty^\pm$, with
$\gamma_\infty=\abs{B_\infty^+-B_\infty^-}/Y_\infty$.
Equations \eqref{eq:absorber-limit} and \eqref{eq:normalized-slab} justify their
limits at $\theta=0$.

In the ideal selective-absorber and reflectionless-plate limit,
\[
 t_+=vQ_\theta\ey,\qquad t_-=vP_yQ_\theta\ey,\qquad r_\pm=0.
\]
Equations \eqref{eq:Qxy}--\eqref{eq:Qyy}, or unitarity of $Q_\theta$, give
\begin{equation}\label{eq:flat-limit}
 Y_0(\omega)=4\left(\frac{b(\omega)}{a(\omega)}\right)^2
          \sin^2\!\left(\frac{\tau\omega a(\omega)}{2a_0}\right),
 \quad C_0^+=Y_0,\quad C_0^-=0.
\end{equation}
Thus $B_0^+=0$, $B_0^-=Y_0$ and the limiting contrast is one.
This is an actual dielectric-stack calculation, not a choice of independent
scattering matrices.

The choices $a\approx a_0$ and
$b\approx a_0/[2\sin(\tau\omega/2)]$ make $Y_0$ arbitrarily close to one
uniformly on $B$. Positivity of the sine on the compact band makes this
approximation stable. For each fixed lossless matrix, first choose finite
$N$ sufficiently large. Next choose positive damping sufficiently small.
The exact finite-layer amplitudes and their normalized Taylor coefficients
depend continuously on that damping. Real-frequency outgoing uniqueness
follows from the passive flux identity and finite-interval ODE uniqueness,
so no real pole obstructs the perturbation. We obtain a fixed, strictly
damped active planar stage whose $Y_\infty(\omega,0)$ is arbitrarily close
to one and whose $\gamma_\infty(\omega,0)$ is arbitrarily close to one,
uniformly on $B$.

\section{Matrix factorization at vanishing coupling}
For the compact aperture, let
\[
 \mathcal H_R=L^2(S_R\times(0,h);\mathbb C)
 \oplus L^2(S_R\times(h,h+d);\mathbb C^2).
\]
Define $J_\theta$ to map the absorber scalar to the $x$ component and the
plate vector through $D_\theta$ into the $xy$ plane. It is a contraction for
$0\le\theta\le1$, and is continuous in norm. It need not be injective at zero.
Set
\[
 \alpha(\omega)=\operatorname{diag}(\alpha_A(\omega),C_\gamma(\omega)/N),
 \quad \chi_\theta=J_\theta\alpha J_\theta^*,\quad
 L_{R,\theta}=\alpha^{-1}-J_\theta^*K_{D_R}(k)J_\theta.
\]
The material weight is
\begin{equation}\label{eq:W}
 W=-\ImH\alpha^{-1}
   =\alpha^{-*}(\ImH\alpha)\alpha^{-1}\succ0.
\end{equation}
For this fixed stage, compactness of $B$ gives a uniform constant $c_*>0$
such that $W\succeq c_*I$. The Green radiation identity implies
\[
 -\Im\ip{r}{L_{R,\theta}r}
 =\ip{r}{Wr}+\Im\ip{J_\theta r}{K_{D_R}J_\theta r}
 \ge c_*\norm r^2.
\]
The coercive sesquilinear form yields invertibility and
$\norm{L_{R,\theta}^{-1}}\le c_*^{-1}$ independently of $R,\omega,\theta$.
This argument does not require $J_\theta$ to be an isometry.

The projected incident field is exactly $J_\theta^*f_\pm=\theta g_\pm$,
where $g_\pm$ is zero in the absorber and $\ey\ee^{\pm ikz}$ in the
two-component plate space. Therefore
\begin{equation}\label{eq:factorization}
 r_{R,\theta}^\pm=L_{R,\theta}^{-1}g_\pm,\quad
 p_\pm=\theta J_\theta r_{R,\theta}^\pm,\quad
 X_R=\theta^2Y_R,\quad A_R^\pm=\theta^2B_R^\pm,
\end{equation}
with
\[
 Y_R=k\Im\ip{g_\pm}{r_{R,\theta}^\pm},\qquad
 B_R^\pm=k\ip{r_{R,\theta}^\pm}{Wr_{R,\theta}^\pm},\qquad
 C_R^\pm=Y_R-B_R^\pm.
\]
For example, with the physical auxiliary current $q=\alpha J_\theta^*E$,
one has $q=\theta r$ and
$E^*(\ImH\chi_\theta)E=q^*Wq$. This also proves the absorption identity.
The normalized quantities extend continuously to zero coupling through
$L^{-1}$, even though the physical incident response then vanishes.

\section{Frequency-uniform finite-aperture localization}
\begin{proposition}\label{prop:aperture}
Freeze positive thicknesses and continuous active matrix laws with
$\ImH\alpha\succ0$ on $B$. On every compact coupling interval in $[0,1]$,
\[
 \sup_{\omega,\theta}
 R^{-1}\norm{r_{R,\theta}^\pm-r_{\infty,\theta}^\pm}_{\mathcal H_R}
 \longrightarrow0,
\]
where the planar current is extended constantly across the aperture. Uniformly
on the same parameters,
\[
 R^{-2}(Y_R,B_R^\pm,C_R^\pm)
 \longrightarrow(Y_\infty,B_\infty^\pm,C_\infty^\pm).
\]
The contrast ratios converge uniformly wherever $Y_\infty$ has a positive
lower bound; here it does on each compact parameter set.
\end{proposition}
\begin{proof}
Let $K_\kappa(k)$ be the transverse Weyl fiber on the fixed height interval,
$\beta=(k^2-\abs\kappa^2)^{1/2}$ with outgoing branch. Its kernels contain
$G_\kappa(z)=i\ee^{i\beta\abs z}/(2\beta)$ and its first derivatives, together
with the longitudinal local term. The explicit bounds from the base paper
are uniform for $k\in[k_1,k_2]$, where $k_i=\omega_i/c_{\rm v}$:
\[
 \norm{K_\kappa}\le C\left(1+\abs\kappa+
              \frac{k^2+\abs\kappa^2}{\abs\beta}\right),\qquad
 \norm{K_\kappa-K_0}\le C_1\abs\kappa
 \quad(\abs\kappa\le k_1/2).
\]
Projection by $J_\theta$ cannot increase these bounds. The low disk is fixed
independently of frequency.

For a fixed smooth window $v$ supported in the unit square and
$v_R(x)=v(x/R)$, the Fourier transform of the localization commutator is
\[
 R^2\widehat v(R\kappa)
 (J_\theta^*K_\kappa J_\theta-J_\theta^*K_0J_\theta)r.
\]
The low disk has global $L^2$ bound $C\norm r\norm{\nabla v}_2$ by Plancherel.
For the complement, Schwartz decay and an integer $m>3$ give
$R^2\abs{\widehat v(R\kappa)}\le C_{v,m}R^{2-m}\abs\kappa^{-m}$.
Its fiber bound is uniformly integrable in $\omega$: on
$k_1/2\le\rho\le2k_2$,
\[
 \abs{k^2-\rho^2}^{-1/2}
 \le k_1^{-1/2}\abs{\rho-k}^{-1/2},
\]
whose radial integral is uniformly finite as the grazing circle moves.
For $\rho>2k_2$, the radial tail is $O(\rho^{2-m})$.
Fourier $L^1$ inversion thus bounds the high-frequency part in transverse
$L^\infty$ by $C_{v,m}R^{2-m}\norm r$, and its aperture $L^2$ norm by
$C_{v,m}R^{3-m}\norm r$. No square of the grazing singularity is used.

For completeness, the transverse zero fiber is
$K_0p(z)=(ik/2)\int\ee^{ik\abs{z-z'}}p(z')\dd z'$. It satisfies
\[
 \Im\ip p{K_0p}=\frac{k}{4}\left(
 \norm{\int\ee^{ikz}p(z)\dd z}^2+
 \norm{\int\ee^{-ikz}p(z)\dd z}^2\right)\ge0.
\]
Thus the same coercivity proof applies to the planar projected operator.
The planar inverse has the same $c_*^{-1}$ bound and its currents have uniformly
bounded height norm. Applying the inverse to the windowed-current residual
gives
\[
 \sup_{\omega,\theta}
 \frac{\norm{r_R-r_\infty}_{\mathcal H_R}}R
 \le C\norm{1-v}_{L^2(S_1)}
       +\frac{C_v}{R}+O(R^{2-m}).
\]
First let $R\to\infty$ for fixed $v$, then approximate one in $L^2(S_1)$.
Density identifies the distributional Fourier construction with the restricted
volume field, as in the base proof. Uniform bounds on $k,W,g$ and currents
of norm $O(R)$ give the work and absorption limits by Cauchy--Schwarz;
energy balance gives the scattering limit.

Finally, the planar kernels, active matrices and drives are continuous in
$(\omega,\theta)$ in operator or Hilbert norm. The inverse identity gives
continuous planar currents. They are nonzero because $g\ne0$ and $L$ is
invertible. Hence $Y_\infty\ge k\ip r{Wr}>0$ pointwise, and continuity on
the compact parameter set supplies a positive minimum. This justifies
uniform division.
\end{proof}

\section{One center-frequency tuning and strict tensor loss}
We now prove the proposed theorem, using a tight bracket so that the floor
holds throughout it. Put $q=E_{\min}/E_*<1$. Choose the oscillator shaping
accuracy so that $Y_0$ is sufficiently close to one that
$\min_B Y_0/Y_0(\omega_c)>q$. Choose finite $N$ and then positive damping as
above, retaining strict material-cap margins, so that the fixed active stage has
\[
 m:=\min_B Y_\infty(\omega,0)>0,\qquad
 \frac{m}{Y_c}>q,\qquad
 \inf_B\gamma_\infty(\omega,0)>1-\epsilon/4,
 \quad Y_c:=Y_\infty(\omega_c,0).
\]
Choose two fixed positive coupling-scale constants $s_-<s_+$ with
\begin{equation}\label{eq:bracket}
 \sqrt{E_{\min}/m}<s_-<\sqrt{E_*/Y_c}<s_+.
\end{equation}
In particular $s_-^2m>E_{\min}$, while the center-frequency extinctions at the
two bracket endpoints straddle $E_*$. The bracket for a physical aperture is
$\theta\in[s_-/R,s_+/R]$.

By uniform continuity and Proposition~\ref{prop:aperture}, for all sufficiently
large finite $R$, this entire bracket lies in $[0,1]$ and satisfies, with
strict margins,
\[
 X_R(\omega,\theta)>E_{\min}\quad(\omega\in B),\qquad
 \Gamma_R(\omega,\theta)>1-3\epsilon/4,
\]
and
\[
 X_R(\omega_c,s_-/R)<E_*<
 X_R(\omega_c,s_+/R).
\]
Indeed, uniformly in $\omega$ and $s\in[s_-,s_+]$,
\[
 R^{-2}Y_R(\omega,s/R)\to Y_\infty(\omega,0),\qquad
 X_R(\omega,s/R)=s^2R^{-2}Y_R(\omega,s/R).
\]
The second relation is exact.
The factor-two bracket from the single-frequency paper would not in general
preserve an arbitrarily high floor; \eqref{eq:bracket} supplies the needed
replacement.

The intermediate tensor is only semidefinitely lossy. For the selected finite
$R$, add the same isotropic passive oscillator $\nu f_0(z)I$ to each layer, where
\[
 f_0(z)=\frac{F_0}{\Omega_0^2-z^2-i\gamma_0z},
 \qquad F_0,\gamma_0>0,\quad\Omega_0>\omega_2.
\]
The full-vector electric inverse before this addition is
\[
 (I-K_{D_R}J_\theta\alpha J_\theta^*)^{-1}
 =I+K_{D_R}J_\theta L_{R,\theta}^{-1}J_\theta^*.
\]
For fixed $R$, $K_{D_R}(k)$ is norm-continuous and uniformly bounded on $B$:
its static singular part is independent of $k$, and its remaining weakly
singular kernel varies continuously on the bounded domain. In detail,
$g_k-g_0=ik/(4\pi)-k^2\abs x/(8\pi)+O(\abs x^2)$ near zero.
The second derivatives of this difference and $k^2g_k$ are $O(\abs x^{-1})$,
uniformly on compact $k$ intervals; an integrable Schur bound gives the asserted
norm continuity. If
$M_R=\sup_B\norm{K_{D_R}(k)}$, the inverse norm is at most $1+M_R/c_*$.
Thus
\[
 \nu\sup_B\abs{f_0}\,M_R(1+M_R/c_*)<1
\]
gives a uniform Neumann perturbation on the band and bracket. Choosing $\nu>0$
still smaller preserves all strict floor, contrast and center-endpoint
margins. Positive real part of $f_0$ preserves the lower real cap; the
preallocated upper margins preserve the upper caps; its positive imaginary
part gives strictly positive tensor loss in every direction.

The final extinction is continuous in $\theta$. A single intermediate-value
argument at $\omega_c$ supplies a fixed $\theta$ in the bracket with
$X(\omega_c)=E_*$ exactly. Since the floor and contrast inequalities hold
throughout the bracket and band, they hold for this one retuned body.
Every selected thickness, aperture, oscillator count and parameter is finite.
The order is essential: shape, choose finite weak/thick stage, give positive
active damping, choose a tight bracket, enlarge the aperture, then add final
isotropic loss and retune. No uniform estimate in all these limits jointly
has been invoked.

Passivity gives $\Gamma\le1$. At any frequency, equality for a compact body
with positive extinction would require one incidence to have zero total
scattering and absorption equal to extinction. Its continuous far field would
then vanish everywhere, contradicting the forward optical theorem. Therefore
$\Gamma(\omega_c)<1$ for every admissible body, so the worst-case supremum
cannot be attained. This completes the analytical construction.

\section{Reproducible planar checks and their limits}
The accompanying \texttt{planar\_checks.py} uses NumPy and exact homogeneous
slab/cascade formulas, with separate standing-wave quadrature of volume loss.
The numerical output is saved in \texttt{numerical\_results.json}.
These tests sample frequency and check algebra; they are not interval
certificates or a three-dimensional finite-aperture calculation.

For a moderate damped tensor, independent volume absorption and outgoing-flux
absorption agree within $2.9\times10^{-16}$; opposite-incidence extinction
agrees at the displayed precision. A separate lossless matrix-plate check
has maximum unitarity residual $9.4\times10^{-16}$ over 31 frequencies.
For the moderate lossless plate matrix
$\left(\begin{smallmatrix}1&0.4\\0.4&1.5\end{smallmatrix}\right)$,
the exact slab plus damped absorber gives the following values over a
21-point grid on $[1,2]$ at $\theta=0.003$:
\[
\begin{array}{c|cc}
N & \max\abs{Y-Y_0} & \min\Gamma\\ \hline
10 & 0.027083 & 0.587407\\
30 & 0.009758 & 0.935600\\
100 & 0.003172 & 0.999431\\
300 & 0.001259 & 0.999437
\end{array}
\]
The $\theta\to0$ limit is essential: increasing $N$ alone at a fixed nonzero
coupling need not drive the contrast all the way to one.

A \emph{separate} rational PSD shaping calculation uses eight poles and
501 frequencies on $[1,2]$, with $\omega_c=1.5$ and target $Y_0=1$.
Its ideal normalized coefficient ranges from $0.997495$ to $1.001077$,
with sampled floor-to-center ratio $0.996766$. The fixed-axis comparison has
ratio $0.75$. The compensation cost is large:
$\max c\approx1.2464\times10^{14}$, while
$\max\abs{a-1}\approx1.04\times10^{-5}$.
This shaped example evaluates \eqref{eq:flat-limit}; it is not the moderate
matrix used in the finite-stack test, and no finite device dimensions are
claimed for it. The analytic argument allows the subsequent scales to be
arbitrarily large but finite. These calculations do not provide useful
engineering resource bounds.

\section{What the draft does and does not settle}
The proposed theorem answers the resource-unrestricted fixed-band question
for every strict floor $0<E_{\min}<E_*$. The tensor shaping is within the
stated two-layer material class, but goes beyond the fixed-principal-axis
rank-one intermediate ansatz used in the original proof.

For that simpler ansatz, a weak plate can give
$q_W=t_W/\ee^{ikd}\to\ee^{i\tau\omega}$ and hence
$Y_0=4\sin^2(\tau\omega/2)$. Center tuning then guarantees only floor ratios
strictly below
\[
 \rho_0=
 \frac{\sin^2(\pi\omega_1/(\omega_1+\omega_2))}
      {\sin^2(\pi\omega_c/(\omega_1+\omega_2))}.
\]
This is a useful restricted-class result, not a bound on the full tensor class.

When $E_{\min}=E_*$, the center must be a global minimum of the exact finite-body
extinction spectrum. Uniform approximation to a flat spectrum and a single
center tuning do not ensure that condition. The present argument establishes
neither existence nor a strict contrast gap in that boundary case.
Exactly constant extinction throughout $B$ is a stronger unresolved condition.
Fixed-resource optima, quantitative resource rates and an all-isotropic
realization also remain outside this construction.

The inherited single-frequency formalization does not verify the new rational
density lemma, matrix aperture extension or broadband selection argument.
Numerical planar calculations can check algebra and illustrate convergence;
they cannot certify the finite three-dimensional Maxwell limit.
No novelty or human-peer-review claim is made.

\begin{thebibliography}{9}
\bibitem{base}
D. Wang and J. Zheng,
\emph{Unit Supremum of Absorption Asymmetry at Prescribed Extinction for
Compact Reciprocal Dielectrics}, manuscript dated 17 October 2026.
Source revision \texttt{22fed75}, Sections 3, 5--9 and Appendix B.
\url{https://dan-k391.github.io/oq7-papers/prescribed-extinction/}.
\bibitem{oscillators}
L. Zhang, F. Monticone and O. D. Miller,
\emph{All electromagnetic scattering bodies are matrix-valued oscillators},
Nature Communications \textbf{14} (2023).
\url{https://doi.org/10.1038/s41467-023-43221-2}.
\end{thebibliography}
\end{document}
